\section{Defenses}
\label{sec:defenses}

\subsection{Turning hooks off is cheaper than assumed}

The obvious mitigation, running the client with hooks disabled, is widely
believed to break too much to be practical. We tested the belief. Installing
the 10{,}000 most downloaded packages of each registry with hooks disabled, in
a clean container, succeeded for 96.2 percent of them. Of the 3.8 percent that
failed, nearly all were native bindings for which a prebuilt artifact exists
and is selected automatically when one is available for the platform. On the
three platforms we tested, the residual failure rate with prebuilt artifacts
enabled was 0.9 percent.

That figure changes the policy conversation. A default of no hooks with an
explicit per-package allowance is a small amount of friction for one package in
a hundred, and it removes the entire class of attack for the other ninety nine.

\subsection{Sandboxing hooks that must run}

For the residual set, the hook must run but does not need ambient authority. A
hook performing native compilation needs the package directory, the toolchain,
and nothing else: no network, no home directory, no environment beyond a
curated set. We ran the residual 0.9 percent under exactly that policy, and all
of them succeeded. The cost is 40~ms of namespace setup per package, which is
noise against a compilation step measured in seconds.

The reason clients do not already do this is not technical difficulty; it is
that the policy has to be the default to matter, and defaults that break one
install in a thousand are hard to ship. Our measurement suggests the true
breakage rate under a compilation-only policy is close to zero.

\subsection{What detection can and cannot add}

Behavioural classification of the kind \tool{} performs is a good registry-side
control and a poor client-side one. Registry side, it runs once per version, at
publication, and a 180-second sandbox execution is affordable at the
publication rates we observe: the busiest of the three registries would need
roughly 40 concurrent sandboxes to keep up. Client side it is useless, because
by the time the client can observe the behaviour the hook has already run.

Detection also inherits the standard limitation. An adversary who reads this
paper will gate the payload on a signal our sandbox does not reproduce, and our
own results show they already gate on continuous integration environment
variables. We treat the 94 percent recall as a measurement of the current
population, not as a durable bound.
